\documentclass[12pt,a4paper]{article}
\usepackage[UTF8,fontset=macnew]{ctex}
\usepackage[margin=2.4cm]{geometry}
\usepackage{booktabs}
\usepackage{tabularx}
\usepackage{setspace}
\setstretch{1.2}
\pagestyle{empty}
\setlength{\parindent}{2em}

\begin{document}

\begin{center}
{\LARGE\bfseries 附件 2}\\[0.6em]
{\Large Cloudra 现阶段学术手稿}\\[0.4em]
{\large Manuscript draft —— 仅供宇树「天才少年」申请备查}
\end{center}

\vspace{1.2em}

\noindent\textbf{状态：}尚未投稿，不是预印本，不是已发表论文。\\
\textbf{请勿：}转发、上网、当作已发表成果引用。\\
\textbf{署名：}Cloudra Team\\
\textbf{申请联系：}杨程驿，\texttt{951899550@qq.com}

\vspace{0.8em}
\section*{阅读时请区分两层}

\renewcommand{\arraystretch}{1.35}
\begin{tabularx}{\textwidth}{X X}
\toprule
\textbf{手稿里已经成立的} & \textbf{手稿里尚未跑完的} \\
\midrule
问题定义（坐标幻觉、「靠嘴宣布完成」） &
100 反应四档消融表中的 Placeholder 数值 \\
系统架构：Delivery DAG、阶段锁、coord\_guard、Runtime v2、督导、GraphRAG 设计 &
三个真实案例正文中的 empirical benchmarking is underway \\
评估协议：四档难度、两阶段消融、早拒答、Progress-Cost &
摘要/结论里按「将要测得」口吻写下的 $0.0\%$ 等句——那是门禁的设计目标，不是已统计结果 \\
工程上已跑通的 Diels--Alder 端到端（见申请书） &
把右栏变成可复查记录，即本次赞助要做的事 \\
\bottomrule
\end{tabularx}

\vspace{1.2em}
\noindent 源文件：申请人仓库 \texttt{paper/main.tex}。本页之后为手稿正文。

\end{document}
